% Pollis drill-down: Compose > Simulation Games, cases played by teams in the chat, round by round,
% on a Julia model run on the facilitator's computer: hosting and joining, teams and roles, a round, news,
% negotiations, the market of a case that trades by the clock, the debrief, the five cases, how a case is written,
% and how the games are made (17 November 2026). For users who want the detail, and for whoever maintains Pollis.
% Compile with: latexmk -pdf workflow_simulation_games.tex
% Shared look for the one-page Pollis workflow handouts (same fonts and colours as the brochure).
\documentclass[10pt,a4paper]{article}

\usepackage[a4paper,left=17mm,right=17mm,top=14mm,bottom=16mm,footskip=8mm]{geometry}
\usepackage[T1]{fontenc}
\usepackage[default,semibold]{sourcesanspro}
\usepackage{microtype}
\usepackage{graphicx}
\usepackage[table]{xcolor}
\usepackage{tikz}
\usepackage[most]{tcolorbox}
\usepackage{fontawesome5}
\usepackage{enumitem}
\usepackage{fancyhdr}
\usepackage[hidelinks]{hyperref}

\graphicspath{{images/}}

% ---------------------------------------------------------------- colours (as in the brochure)
\definecolor{navy}{HTML}{1C3A5E}
\definecolor{pblue}{HTML}{3B78B5}
\definecolor{amber}{HTML}{F2B33D}
\definecolor{deepamber}{HTML}{D9902A}
\definecolor{pale}{HTML}{EEF3F9}
\definecolor{ink}{HTML}{2A2F36}
\definecolor{logonavy}{HTML}{173D63}
\definecolor{logodark}{HTML}{C98A1E}
\definecolor{logoamber}{HTML}{F2B446}
\definecolor{logocream}{HTML}{FEEFD2}
\definecolor{soft}{HTML}{5E6773}
\color{ink}

\setlength{\parindent}{0pt}
\setlength{\parskip}{3pt}
\setlist{nosep,leftmargin=*,itemsep=2pt}

\pagestyle{fancy}
\fancyhf{}
\renewcommand{\headrulewidth}{0pt}
\fancyfoot[L]{\footnotesize\color{soft}\textbf{Pollis}\enspace Bright Insights through Exploration, Modelling, Simulation and Optimisation}
\fancyfoot[R]{\footnotesize\color{soft}github.com/saragga/pollis/releases/latest}

% The Pollis logo, redrawn so it prints sharply at any size
\newcommand{\pollislogo}[1]{%
  \begin{tikzpicture}[x=#1/512,y=-#1/512,baseline=(current bounding box.center)]
    \fill[logonavy] (0,0) rectangle (512,512);
    \fill[logodark] (36,422) -- (104,256) -- (172,422) -- cycle;
    \fill[logoamber] (104,422) -- (203,180) -- (302,422) -- cycle;
    \fill[logocream] (203,422) -- (339,89) -- (476,422) -- cycle;
  \end{tikzpicture}}

% Top of the page: logo, name, kind of document (Workflow unless given, e.g. [Drill-Down]) and its subject;
% then the question and a one-line promise
\newcommand{\workflowhead}[4][Workflow]{%
  \thispagestyle{fancy}%
  \begin{minipage}[c]{13mm}\pollislogo{11mm}\end{minipage}%
  \begin{minipage}[c]{\dimexpr\linewidth-13mm\relax}
    {\color{navy}\Large\bfseries Pollis}\enspace{\color{soft}#1}\hfill
    {\color{deepamber}\small\bfseries\MakeUppercase{#2}}
  \end{minipage}\par\vspace{7mm}
  {\color{navy}\fontsize{22}{26}\selectfont\bfseries #3}\par\vspace{4pt}
  {\color{soft}\large #4}\par\vspace{5mm}}

% A menu path: \menu{Explore}{Access Data Libraries}... is written as \path{Explore \gt Access Data Libraries}
\newcommand{\gt}{\,{\color{deepamber}\bfseries ›}\,}
\newcommand{\menupath}[1]{{\footnotesize\bfseries\color{navy}#1}}

% A screenshot if images/<file> exists, otherwise an orange placeholder that says what to capture
\newlength{\shotheight}\setlength{\shotheight}{43mm}
\newcommand{\stepshot}[2]{%
  \IfFileExists{images/#1}{%
    \tcbox[enhanced, colback=black, colframe=black!70, boxrule=0.4pt, arc=2pt,
      left=0pt,right=0pt,top=0pt,bottom=0pt, boxsep=0pt, halign=center]{%
      \includegraphics[width=\dimexpr\linewidth-1pt\relax,height=\shotheight,keepaspectratio]{#1}}}{%
    \begin{tcolorbox}[enhanced, colback=white, colframe=orange!80, boxrule=0.6pt, arc=2pt,
      height=\shotheight, valign=center, halign=center, left=6pt, right=6pt]
      {\color{orange}\Large\faCamera}\par\smallskip
      {\color{orange}\footnotesize\bfseries Screenshot: #2}\par\smallskip
      {\color{orange!80!black}\scriptsize images/#1}
    \end{tcolorbox}}}

% One step: number, title, menu path, image file, what to capture, what it means
\newcommand{\step}[6]{%
  \begin{tcolorbox}
    {\color{deepamber}\Large\bfseries #1}\enspace{\color{navy}\large\bfseries #2}\par\vspace{1pt}
    \menupath{#3}\par\vspace{4pt}
    \stepshot{#4}{#5}\par\vspace{3pt}
    {\small #6}
  \end{tcolorbox}}

\newenvironment{steps}{%
  \begin{tcbraster}[raster columns=2, raster equal height=rows, raster column skip=4mm,
    raster row skip=4mm, enhanced, colback=pale, colframe=pale, boxrule=0pt, sharp corners,
    left=6pt,right=6pt,top=6pt,bottom=6pt, fontupper=\raggedright]}{\end{tcbraster}}

% Boxes for the bottom of the page
\newtcolorbox{pbox}[2][]{enhanced, colback=pale, colframe=amber, boxrule=0pt,
  leftrule=3pt, sharp corners, left=7pt, right=7pt, top=6pt, bottom=6pt,
  fontupper=\small\raggedright, before upper={{\color{navy}\normalsize\bfseries #2}\par\smallskip}, #1}
\newtcolorbox{nbox}[1][]{enhanced, colback=navy, colframe=navy, boxrule=0pt,
  sharp corners, left=8pt, right=8pt, top=6pt, bottom=6pt, coltext=white,
  fontupper=\small\raggedright, #1}

% A Next Steps item inside the bottom boxes
\newcommand{\nsitem}[2]{\item[\color{deepamber}\faChevronRight] \textbf{#1}: #2}

% Citations: \dcite{key}{Breiman (2001)} links to its entry \dref{key}{...} in the References, keyed as in pollis.bib
\newcommand{\dcite}[2]{\hyperlink{ref:#1}{\color{pblue}#2}}
\newcommand{\dref}[2]{\raisebox{1.1em}[0pt][0pt]{\hypertarget{ref:#1}{}}#2\par}
\newcommand{\ddoi}[1]{\href{https://doi.org/#1}{\color{pblue}doi:#1}}
\newcommand{\references}{\section*{References}\addcontentsline{toc}{section}{References}\footnotesize\raggedright\setlength{\parskip}{3pt}}

% The page of each heading, for Ask Pollis
% Pollis: the page of each heading, for Ask Pollis. A document that inputs this file, at the end of its preamble, writes
% <jobname>.labels beside its PDF, a line page|level|heading for each heading, the page counted from the first, the
% cover, as 1. compile_documents.jl copies it to ~/.pollis/documents, and ask_index.jl cuts the document into passages
% under its headings with it. \asklabel{1}{Title} marks a topic or a page, \asklabel{2}{Title} a heading inside it; the
% sections and subsections of a document that has them are marked as they go into its contents. A package vignette
% inputs it before its \documentclass (compile_vignette.jl), so the contents are taken over when the document begins,
% after the packages have taken them over. Read a second time, as by a copy of a document for Edit the LaTeX Source,
% which inputs it as resources/ask_labels when Pollis has put it in already (document_source.jl), it does nothing.
\ifdefined\asklabel\endinput\fi
\makeatletter
\newwrite\pollis@labels
\immediate\openout\pollis@labels=\jobname.labels
\newcommand{\asklabel}[2]{\write\pollis@labels{\the\ReadonlyShipoutCounter|#1|\unexpanded{#2}}}
\def\pollis@toc{toc}
\newcommand{\pollis@contentsline}[3]{%
  \def\pollis@file{#1}%
  \ifx\pollis@file\pollis@toc\asklabel{#2}{#3}\fi
  \pollis@addcontentsline{#1}{#2}{#3}}
% The headings of the New Methods style of the vignettes, \nmsection and \nmsub, which go into no contents
\def\pollis@asksection#1{\pollis@nmsection{#1}\asklabel{1}{#1}}
\def\pollis@asksub#1{\pollis@nmsub{#1}\asklabel{2}{#1}}
\AddToHook{begindocument/end}{%
  \let\pollis@addcontentsline\addcontentsline\let\addcontentsline\pollis@contentsline
  \ifdefined\nmsection\let\pollis@nmsection\nmsection\let\nmsection\pollis@asksection\fi
  \ifdefined\nmsub\let\pollis@nmsub\nmsub\let\nmsub\pollis@asksub\fi}
\makeatother



\usepackage{titlesec}
\usepackage{tabularx}
\usepackage{booktabs}
\usepackage{listings}

\fancyhf{}
\fancyfoot[L]{\footnotesize\color{soft}\textbf{Pollis}\enspace Drill-down: Simulation Games}
\fancyfoot[R]{\footnotesize\color{soft}\thepage}

\titleformat{\section}{\color{navy}\large\bfseries}{\color{deepamber}\thesection}{0.6em}{}
\titleformat{\subsection}{\color{navy}\normalsize\bfseries}{\color{deepamber}\thesubsection}{0.5em}{}
\titlespacing*{\section}{0pt}{12pt}{4pt}
\titlespacing*{\subsection}{0pt}{8pt}{2pt}
\setlength{\parskip}{4pt}

\lstset{basicstyle=\ttfamily\footnotesize, columns=fullflexible, keepspaces=true, breaklines=true,
  backgroundcolor=\color{pale}, frame=leftline, rulecolor=\color{amber}, framerule=2pt, xleftmargin=6pt,
  framexleftmargin=6pt, aboveskip=4pt, belowskip=4pt, upquote=true, tabsize=2}

\usepackage{xurl}
\newcommand{\file}[1]{{\small\nolinkurl{#1}}}
\newcolumntype{L}{>{\raggedright\arraybackslash}X}
\newcommand{\ui}[1]{\textbf{#1}}
\newcommand{\setting}[1]{\file{#1}}
% Code with spaces in it, such as a command line, which \file would run together
\newcommand{\code}[1]{\texttt{\small #1}}
% A menu path at the start of a table cell
\newcommand{\cellpath}[1]{\leavevmode\menupath{#1}}

\begin{document}

\workflowhead[Drill-Down]{Simulation Games}
  {Teams, rounds and a model the class can open}
  {How a class plays a case in the chat, team against team, round by round, on a Julia model run on the
   facilitator's computer, and how the debrief turns every decision, result and message into tables.}

\begin{nbox}
{\color{amber}\bfseries Drill-down.}\enspace
A look under the surface of \textbf{Compose \gt Simulation Games}. It is for anyone who wants to
run a game with a class, know what each player sees and what is kept, or write a case of their own;
\S\ref{sec:machinery} is for whoever maintains Pollis. It describes the state of 17 November 2026.
\end{nbox}

\tableofcontents

\section{What Simulation Games are}\label{sec:what}
In a simulation game the players run something, a company, a deal, a trading desk, in a world a model keeps: they
decide, the model answers with results, and they decide again in its light. Business games of this kind have been
used in teaching for decades; \dcite{FHWG2009a}{Faria, Hutchinson, Wellington and Gold (2009)} review how they
developed over forty years, in the technology they run on, why they are adopted and how they are administered.

In Pollis, a game is a session of the Collaboration chat. A \emph{facilitator}, the teacher, hosts it; the players
join it with its invitation and are split into teams. Each round, every team submits its decisions in a form beside
the chat, and the facilitator's copy of Pollis runs the case's model, in Julia, on the facilitator's computer. The
model is not a black box: it is a file of Julia code that comes with the case, which the class opens in the
debrief.

\begin{center}\small
\renewcommand{\arraystretch}{1.25}
\begin{tabularx}{\linewidth}{@{}p{0.42\linewidth}L@{}}
\toprule
\textbf{Menu item} & \textbf{What it is for}\\
\midrule
\cellpath{Simulation Games \gt Host a Game...} & The facilitator: chooses a case and starts the game's session\\
\cellpath{Simulation Games \gt Join a Game...} & A player: joins with the invitation the facilitator sent\\
\cellpath{Simulation Games \gt Next Round} & The facilitator: forms the teams, then plays each round\\
\cellpath{Simulation Games \gt Release News...} & The facilitator: releases an item of news while a round is decided\\
\cellpath{Simulation Games \gt Open the Debrief} & The facilitator: writes the debrief of a game, during it or after\\
\bottomrule
\end{tabularx}
\end{center}
Next Round, Release News... and Open the Debrief are also buttons in the facilitator's column of the chat page.

\section{A game from start to end}\label{sec:game}
\begin{enumerate}
  \item \textbf{Host.} \ui{Host a Game...} lists the cases of the installed toolboxes, by their toolbox's name, and
    \ui{Choose a Case Folder...} for a case of your own. When the case has several ways of organising its teams,
    it asks which; when it can be played in different lengths, it asks how many rounds. The session starts, and
    its invitation ends in \file{@game}.
  \item \textbf{Join.} Each player chooses \ui{Join a Game...}, pastes the invitation and types their name. A game
    is a class session, so it takes up to 30 players, each connected to the facilitator only.
  \item \textbf{Teams.} The first \ui{Next Round} asks for the number of teams, within the case's bounds and at most
    one per player, splits the players at random, and lets the facilitator move any of them before starting.
    Players who join later are placed by \menupath{Classroom \gt Split into Breakout Groups...}, which keeps the teams
    and lets the facilitator move anyone.
  \item \textbf{Start.} The model's \file{initial} sets up the world; what it says before the first round is shown
    as the opening, and each player reads the case's briefing and that of their role.
  \item \textbf{Rounds.} Each team submits its decisions; the facilitator sees which teams have, and presses
    \ui{Next Round}, which asks first when some have not. The model plays the round, and everyone sees the
    results. The facilitator can release news while a round is decided; the teams can write to each other.
  \item \textbf{End.} The game ends after its last round, or earlier when the model says so, such as when a plan is
    agreed, or when fewer than two teams remain. \ui{Open the Debrief} then takes the place of Next Round.
\end{enumerate}

\section{Teams, roles and briefings}\label{sec:roles}
A case can have several \emph{structures}, ways of organising its teams, and the facilitator chooses one when
hosting. Each structure has \emph{roles}, each with its briefing and the decisions it sets. In Grow a Chocolate
Chain, for example, \ui{Board Members} has one role, which reads every briefing and sets every decision, while
\ui{Departments} has four: the Chief Executive sets the shops and the staff, Marketing the price and the marketing,
Sourcing the beans, and Finance nothing, but reads the costs and the cash, so that the team has to talk.

\begin{itemize}
  \item \textbf{Dealing the roles.} Roles are dealt when the game starts and whenever a team changes. A player keeps
    their roles while they stay in their team; a newcomer takes a role nobody in the team holds, or else shares
    the one fewest hold; a role still unheld goes to the player with fewest. When a player leaves, their roles go
    to the others.
  \item \textbf{Submitting.} A player's submission changes only the decisions their roles set, and the last one
    submitted counts. The fields of the other roles are shown, but cannot be changed (``Set by Marketing''). The
    facilitator sees who submitted, and which roles have not.
  \item \textbf{Briefings.} The case's briefing, which everyone reads, and the briefings of the player's roles are
    beside the chat, in sections that can be closed. The facilitator reads every briefing.
\end{itemize}

\subsection{Sides}
In some cases the teams are not alike: a buyer and the sellers, or shareholders, banks and bondholders. Each kind
of team is a \emph{side}, with its own number of teams, briefing, roles, decisions and private results. In such a
case the debrief ranks the teams of each side among themselves; a side of one team is set against doing nothing.

\section{A round}\label{sec:round}
\begin{itemize}
  \item \textbf{The decisions.} Each is a number, with its unit, bounds and step, or a choice among options. A
    decision can be made only by some sides, only in some rounds, or once for each other team, such as a bid for
    each rival. A team with nothing to decide in a round is told so.
  \item \textbf{The model.} \ui{Next Round} writes every team's decisions and runs the case's model in a fresh Julia
    process on the facilitator's computer. Its random numbers come from a seed set when the game starts and the
    round's number, so a round played again gives the same results. The Julia run is the one the setting
    \setting{julia.executablePath} names, else the \file{julia} a terminal would find; the players need no Julia.
  \item \textbf{The results.} The model returns the public results, a row for each team that everyone sees, and the
    private ones, which each team alone sees, with a text for everyone and, when it has one, a text for each team.
    It can take a team out of the game, a company bought, say, with the reason; the team keeps its results, but
    decides no more.
\end{itemize}

\section{News}\label{sec:news}
A case brings items of news, each with the first round in which it can be released, a title, a text, the roles who
read it, and the values of the model's parameters it changes. An item can also leave some of this to chance: a
\file{chance} of being in the game at all, a window of rounds, \file{rounds = [2, 4]}, in place of one round, and
a change given as a range, \file{[-0.4, -0.2]}, in place of one value. When the game starts, Pollis draws each
from the game's seed, so that a student who has read \file{case.toml} knows what may happen, not what will.
\ui{Release News...} offers the items drawn for the round, and earlier ones not yet released, or \ui{Write a News Item...}: a title, a text, any parameters changed, picked
from the case's numbers with their values now, and the roles who read it. The players are not told which
parameters changed; the facilitator's card lists each change, from and to. News is released only while a round is
decided, and its changes apply from that round on. A player sees an item as a card in the chat,
``News for Sourcing, Finance: \dots'' when it is not for everyone; latecomers see the news released so far.

\section{Negotiations}\label{sec:negotiations}
Unless a case says otherwise, a player can write to another team while a round is decided: \ui{Write to}, the team,
a text, and \ui{Send to the Team}. The message goes to everyone in both teams, and to the facilitator, who reads it;
the players are told so in the form and when they join a team. Nothing makes a team keep what it promised. The chat
shows ``Ana to Team 2'' to the writer's team, ``Ana (Team 1) to your team'' to the other, and ``Ana (Team 1) to
Team 2'' to the facilitator, and every message is kept for the debrief.

\section{A market by the clock}\label{sec:market}
A case with a market, such as Trading Floor, is played in \emph{periods} instead of rounds of decisions: each period,
the teams are trading desks that buy and sell one asset on an order book while a clock runs, a few minutes, beside
the chat. Classroom markets have a long history: \dcite{Smi1962a}{Smith (1962)} had students trade as buyers and
sellers in a double auction, each with a value or a cost only they knew, and found prices close to the competitive
equilibrium; \dcite{GS1993a}{Gode and Sunder (1993)} found the allocative efficiency of a double auction to derive
largely from its structure, even with traders who submit random bids and offers within their budgets.

\subsection{The rules of the book}
\begin{itemize}
  \item An order buys or sells a number of shares at a price or better, the price in steps of the case's tick.
  \item An order that can trade at once does: with the best price on the other side first and, at the same price,
    the oldest order first, at the price of the order already on the book.
  \item What does not trade stays on the book until another desk trades with it, someone in the desk cancels it,
    or the period closes. Orders still on the book at the close are cancelled.
  \item A desk can buy only with the cash its other buy orders do not hold, and sell only the shares its other sell
    orders do not hold: no borrowing and no selling short. An order that would trade with the desk's own order is
    refused. A desk has at most 20 orders on the book.
\end{itemize}
The page shows the asset, the last trade, the volume, the time left, the best eight prices to buy and to sell with
their shares, and the last twelve trades; a player also sees their desk's cash and shares and their own orders,
with \ui{Cancel}. The facilitator sees every desk.

\subsection{The close of a period}
The period closes when its time is up; \ui{Next Round} closes it earlier, after asking. Next Round then announces
the period's results: each desk's position, its cash, shares, what it bought and sold and at what average prices,
and the market's last, average, highest and lowest price and its volume, go to the model as the desk's decisions.
Every trade is kept, for the debrief.

\subsection{Private information}
In Trading Floor some desks receive a hint of the company's results before a period. Whether a market's prices
reveal what some traders know is the question of \dcite{PS1988a}{Plott and Sunder (1988)}, who found that laboratory
markets aggregated the traders' diverse information well in some designs, and hardly at all in others. The debrief
lets the class ask the same of their own market: did the price move toward the value before the results were
announced?

\section{The debrief}\label{sec:debrief}
The debrief, the analysis that follows the game, is where a simulation teaches: \dcite{Led1992a}{Lederman (1992)}
calls it the heart of experiential learning, and reviews its phases and how to conduct it. \ui{Open the Debrief}
works on the game hosted, once a round has been played, or, outside a game, on a game folder chosen. It writes a
folder \file{debrief} in the game's folder, replacing an earlier one, and opens \file{debrief.md} in the Markdown
preview.

\begin{center}\small
\renewcommand{\arraystretch}{1.2}
\begin{tabularx}{\linewidth}{@{}p{0.27\linewidth}L@{}}
\toprule
\textbf{File} & \textbf{What it holds}\\
\midrule
\file{debrief.md} & The teams and their structure; the standings after the last round played, by the case's
  score, side by side in a case with sides; each round's text, public results, news with its readers and changes,
  the number of messages, and Behind the Results, what the model wrote for the debrief alone; every message between teams; the case's questions for the class\\
\file{decisions.csv} & Every team's decisions in every round\\
\file{results.csv} & The public results of every round\\
\file{team_results.csv} & The private results of every team and round\\
\file{news.csv} & The news released, with its round, readers and changes\\
\file{negotiations.csv} & Every message, with its time, writer and teams\\
\file{trades.csv} & In a case with a market: every trade, with its period, time, price and shares, the two desks, and who placed each
  order\\
\file{case/} & A copy of the case, so that the class can open \file{model.jl}\\
\bottomrule
\end{tabularx}
\end{center}
In a live game the facilitator is asked whether to share the debrief with the players, and warned that they will
see every team's accounts and messages. Shared, \file{debrief.md}, \file{model.jl} and \file{case.toml} go to the
players as files, and the tables as datasets, to open with the ordinary panels.

\section{The five cases}\label{sec:cases}
Each case comes in a toolbox of its own, installed from the Toolboxes section of the Extensions view. The companies
and their names are made up.

\begin{center}\small
\renewcommand{\arraystretch}{1.25}
\begin{tabularx}{\linewidth}{@{}>{\raggedright\arraybackslash}p{0.19\linewidth}>{\raggedright\arraybackslash}p{0.23\linewidth}L@{}}
\toprule
\textbf{Case} & \textbf{Teams and rounds} & \textbf{What the teams do}\\
\midrule
Grow a Chocolate Chain & 2 to 6 brands, 5 years & Each runs a chocolate brand in one country, selling in its own
  shops: its price, beans, shops, staff and marketing. The highest cumulative profit wins\\
Merge or Be Bought & 3 bicycle makers, 6 quarters & Each bids for its rivals' shares and sets the least price it
  accepts for its own; a company bought leaves the game. The best return for the original shareholders wins\\
A Vineyard for the Portfolio & A drinks group and 3 wineries, 6 rounds & The buyer looks into the wineries,
  offers cash and an earn-out, and runs what it bought for three years; each winery's owners know a problem the
  buyer does not. The value created, on each side\\
Save the Company & Shareholders, banks and bondholders of a retailer, 4 or 6 quarters & The shareholders propose a
  restructuring; the banks and the bondholders accept it, or call in the loan, or sell. A plan agreed, or the
  administration. The value at the end, on each side\\
Trading Floor & 3 to 8 desks, 4 or 6 periods & The desks trade one company's shares on the order book while its
  results move their value; some desks receive hints first. The best return wins\\
\bottomrule
\end{tabularx}
\end{center}

\subsection{What each case is about}
\begin{itemize}
  \item \textbf{Merge or Be Bought.} A bid is made only when the bidder's cash and borrowing cover it, and a company
    is bought when its highest bid reaches its owners' least price. The synergy a buyer expects is its own, and
    known to it alone. \dcite{Rol1986a}{Roll (1986)} argued that bidders who trust their own valuation too much simply
    pay too much for their targets; the debrief asks whether the class's buyers did.
  \item \textbf{A Vineyard for the Portfolio.} Sellers who know more than the buyer are the setting of
    \dcite{Ake1970a}{Akerlof (1970)}, where uncertainty about quality drives good products from the market. The
    buyer can pay to find out, or protect itself with an earn-out, paid when the winery reaches its plan:
    \dcite{KA2000a}{Kohers and Ang (2000)} found earnouts used against the misvaluation of targets about which
    information is asymmetric, and as a bonus that keeps the target's managers.
  \item \textbf{Save the Company.} The banks are secured and paid first, the bondholders after them, the
    shareholders last. A restructuring agreed outside the courts can leave every side better off than
    administration, if they agree in time: \dcite{GJL1990a}{Gilson, John and Lang (1990)} studied firms in default
    that restructured their debt privately, and those that went to court instead.
  \item \textbf{Trading Floor.} The company's value moves at each close by its results, and by the news the
    facilitator releases; at the end the company is taken over at its value, and every desk's shares are paid in
    cash. Before each period, every other desk receives its analysts' hint of the results (\S\ref{sec:market}).
\end{itemize}

\section{Writing a case}\label{sec:writing}
A case is a folder with \file{case.toml}, which says what the players see and decide, and \file{model.jl}, which
says what happens; its briefings are Markdown files beside them.

\subsection{Create a Simulation Game}
\ui{Compose > Simulation Games > Create a Simulation Game...} writes a case in four steps.
\begin{enumerate}
  \item \textbf{Start From.} A case to start from: \ui{Sell One Product}, the minimal case of Pollis (one price,
    three rounds), or a case of the installed toolboxes as a pattern: rivals in one market (Grow a Chocolate Chain),
    bids among equals (Merge or Be Bought), a buyer and sellers with secrets (A Vineyard for the Portfolio), a
    negotiation over a plan (Save the Company) or a market by the clock (Trading Floor). \ui{Create the Case}
    copies it under the title given into \file{games/cases} of the Pollis Workspace, else
    \file{~/.pollis/games/cases}. \ui{Your Cases} lists those cases, and any opened with
    \ui{Open a Case Folder...}
  \item \textbf{Write.} The case's files, this drill-down, and \ui{Write with Agent}: the coding agent of
    \ui{Compose > Coding Agents}, with the skill \file{write-simulation-game}, proposes a design from what the user
    describes, writes the files once the user agrees, and plays every round with \file{game.jl} until none fails.
  \item \textbf{Test.} \ui{Test the Case} reads \file{case.toml} as Host a Game does, plays every round with the
    fewest teams the case allows (at least two), each on its default decisions, and opens the debrief of the test
    game. It lists what the results' rows lack of the columns of \file{case.toml}, and what they have that no player
    sees; when a round fails, it shows the model's error, and \ui{Fix with Agent} hands it to the agent.
  \item \textbf{Share.} \ui{Host This Case} opens Host a Game with the case chosen; \ui{Show the Folder}, to give it
    to colleagues; \ui{Package with Agent}, to make it a toolbox of cases (below).
\end{enumerate}
What the players must not know is best drawn for each game: anyone can read \file{model.jl}, so a secret fixed in
it is no secret. Give it as a range \code{[low, high]} in \file{[parameters]}, draw it in \code{initial}, give the
range in the briefings, and write what was drawn in the \file{debrief} of \code{initial}, as the five cases do.

\subsection{case.toml}
\begin{center}\small
\renewcommand{\arraystretch}{1.2}
\begin{tabularx}{\linewidth}{@{}p{0.27\linewidth}L@{}}
\toprule
\textbf{Key} & \textbf{What it says}\\
\midrule
\file{title}, \file{description} & The case's name and what it is about\\
\file{rounds}, \file{round}, \file{round_names} & The number of rounds, the word for one (\ui{Year}, \ui{Quarter},
  \ui{Period}) and, when they differ, each round's name\\
\file{round_choices} & The lengths the facilitator chooses from, such as \code{[4, 6]}\\
\file{teams}, \file{team_names} & The least and most teams, and their names\\
\file{briefing}, \file{debrief} & The briefing everyone reads, and the questions for the debrief\\
\file{negotiations} & \code{false} when the teams cannot write to each other\\
\file{[[sides]]} & The kinds of team: key, label, number of teams and briefing\\
\file{[[structures]]} & The ways the teams can be organised, each with \file{[[structures.roles]]}: key, label,
  briefing and \file{decides}, the decisions the role sets\\
\file{[[decisions]]} & Key, label, \file{kind} (\code{number} or \code{choice}), unit, bounds, step, options,
  default and help; \file{rivals}, \file{sides} and \file{rounds} when a decision is not made by every team in every
  round\\
\file{[results]} & \file{winner}, the column that ranks the teams, and the columns of \file{[[results.public]]} and
  \file{[[results.private]]}, each with key, label, unit and digits\\
\file{[[news]]} & Round, or a window of \file{rounds}, a \file{chance}, title, text, roles, and
  \file{[news.parameters]}, the parameters it changes, each to a value or a range \file{[low, high]}\\
\file{[market]} & For a case that trades by the clock: the asset, the currency, the minutes of a period, the tick,
  each desk's cash and shares, and the largest order and price\\
\file{[parameters]} & The model's numbers, which news can change\\
\bottomrule
\end{tabularx}
\end{center}

\subsection{model.jl}
The model defines two functions. Pollis loads it in a module of its own, with Random, Statistics, LinearAlgebra and
Printf; other packages the case needs are listed in \file{packages}, and installed when missing.
\begin{lstlisting}
function initial(rng, parameters, teams)
    # the world before round 1: return state, or (state, results) to show an opening
end

function play(rng, state, decisions, parameters)
    # one round: return (state, results)
end
\end{lstlisting}
\begin{itemize}
  \item \file{parameters} is \file{[parameters]}, as the news released so far changed it; in a case with
    \file{round_choices} it holds \file{rounds}, the length chosen, and in a case with a market, \file{market}.
  \item \file{teams} and \file{decisions} have, for each team, its number, its name and, in a case with sides, its
    side; \file{decisions} also the values the team chose, by key, and a decision made for each rival is a
    dictionary by the rival's number.
  \item \file{state} is a dictionary of what TOML can hold, written to the round's folder, so that a game can be
    read and played again.
  \item \file{results} has \file{public}, \file{private} by team, and when needed \file{text}, \file{private_text},
    \file{out} (the teams that leave, with the reason), \file{over} and \file{debrief}, a Markdown text no player
    sees during the game, which the debrief shows under the round: the secrets the model drew for the game, say, in
    the results of \file{initial}, or how far apart two teams' prices were in a negotiation.
\end{itemize}

\subsection{A toolbox of cases}
A toolbox brings cases by listing them in its \file{package.json}, under \file{pollisToolboxes}, as \file{games},
each the \file{folder} of a case; a toolbox of cases alone has no panels. Installed, its cases are in
\ui{Host a Game...}. To try a case before packaging it, choose \ui{Choose a Case Folder...}. \ui{Package with
Agent}, in the Share step of Create a Simulation Game, writes such a toolbox with the user (its \file{package.json},
README and the case in \file{games/}) and packs it as a \file{.vsix} file, which others install with
\ui{Extensions: Install from VSIX...}

\section{Under the surface}\label{sec:machinery}
\begin{itemize}
  \item \textbf{A game is a class.} The session is a class session (\file{collaborationService.ts}): the facilitator's
    copy joins one room per seat, each player is in one room with the facilitator only, and the facilitator's
    copy passes everything on, to a team or to everyone. Teams are breakout groups that last the game.
  \item \textbf{The messages.} \file{collaborationProtocol.ts} defines them: \file{game} (the round, the team, its
    decisions, roles and briefings), \file{decisions}, \file{gameResults}, \file{negotiation}, and, for a market,
    \file{order}, \file{cancelOrder} and \file{market}, the book as each desk sees it, with its refusals.
  \item \textbf{The run.} \file{GameRun} (\file{collaborationGame.ts}) reads the case (\file{loadCase}), forms the
    teams, deals the roles, writes the round's files and runs \file{media/games/game.jl}, which loads the model,
    applies the news, calls \file{initial} or \file{play} and writes the state and the results as JSON. A round that
    fails leaves \file{error.txt} saying why.
  \item \textbf{The market.} \file{MarketBook} (\file{collaborationMarket.ts}) is kept by the facilitator's copy.
    Prices are kept in whole ticks, so that no rounding creeps into cash, and the cash and shares a desk's orders
    hold are set aside until they trade or are cancelled.
  \item \textbf{The page.} \file{collaborationChat.ts} draws the decisions form, the briefings, the negotiations,
    the market and the facilitator's column, beside the chat.
  \item \textbf{The debrief.} \file{collaborationDebrief.ts} reads a game's folder from disk and writes
    \file{debrief/}.
\end{itemize}

\subsection{A game's folder}
A game is a folder named by the date and time it started, in \file{games/<case>} of the Pollis Workspace, else in
\file{~/.pollis/games/<case>}:
\begin{center}\small
\renewcommand{\arraystretch}{1.2}
\begin{tabularx}{\linewidth}{@{}p{0.32\linewidth}L@{}}
\toprule
\textbf{File} & \textbf{What it holds}\\
\midrule
\file{game.toml} & The case and its folder, the structure, the number of rounds, the seed and the teams\\
\file{round-00/state.toml} & The world before round 1\\
\file{round-k/decisions.toml} & Every team's decisions, with its players and their roles\\
\file{round-k/news.toml} & The news released in the round\\
\file{round-k/negotiations.toml} & The messages between teams\\
\file{round-k/trades.toml} & In a case with a market: the period's trades\\
\file{round-k/state.toml}, \file{results.json} & The world after the round, and its results\\
\bottomrule
\end{tabularx}
\end{center}

\section{Where things are}

\begin{center}\small
\renewcommand{\arraystretch}{1.2}
\begin{tabularx}{\linewidth}{@{}p{0.46\linewidth}L@{}}
\toprule
\textbf{File} & \textbf{Role}\\
\midrule
\file{.../collaboration/collaborationGame.ts} & Cases, teams, roles and the rounds\\
\file{.../collaboration/collaborationMarket.ts} & The order book\\
\file{.../collaboration/collaborationDebrief.ts} & The debrief\\
\file{.../collaboration/collaborationService.ts} & The session, the menu's commands\\
\file{.../collaboration/collaborationChat.ts} & The chat page\\
\file{.../collaboration/collaborationProtocol.ts} & The messages\\
\file{.../collaboration/createGamePage.ts} & Create a Simulation Game\\
\file{.../compose/browser/media/games/game.jl} & The runner of a case's model\\
\file{.../compose/browser/media/games/minimal/} & Sell One Product, the minimal case\\
\file{write-simulation-game} (Pollis Agents plugin) & The skill of Write, Fix and Package with Agent\\
\file{pollis-toolboxes}, \file{pollis-toolbox-<case>} & The five cases\\
\bottomrule
\end{tabularx}
\end{center}

\references
\dref{Ake1970a}{Akerlof, G.\,A. (1970). The Market for ``Lemons'': Quality Uncertainty and the Market Mechanism. \emph{The Quarterly Journal of Economics} 84(3), 488--500. \ddoi{10.2307/1879431}}
\dref{FHWG2009a}{Faria, A.\,J., Hutchinson, D., Wellington, W.\,J. and Gold, S. (2009). Developments in Business Gaming: A Review of the Past 40 Years. \emph{Simulation \& Gaming} 40(4), 464--487. \ddoi{10.1177/1046878108327585}}
\dref{GJL1990a}{Gilson, S.\,C., John, K. and Lang, L.\,H.\,P. (1990). Troubled Debt Restructurings: An Empirical Study of Private Reorganization of Firms in Default. \emph{Journal of Financial Economics} 27(2), 315--353. \ddoi{10.1016/0304-405X(90)90059-9}}
\dref{GS1993a}{Gode, D.\,K. and Sunder, S. (1993). Allocative Efficiency of Markets with Zero-Intelligence Traders: Market as a Partial Substitute for Individual Rationality. \emph{Journal of Political Economy} 101(1), 119--137. \ddoi{10.1086/261868}}
\dref{KA2000a}{Kohers, N. and Ang, J. (2000). Earnouts in Mergers: Agreeing to Disagree and Agreeing to Stay. \emph{The Journal of Business} 73(3), 445--476. \ddoi{10.1086/209649}}
\dref{Led1992a}{Lederman, L.\,C. (1992). Debriefing: Toward a Systematic Assessment of Theory and Practice. \emph{Simulation \& Gaming} 23(2), 145--160. \ddoi{10.1177/1046878192232003}}
\dref{PS1988a}{Plott, C.\,R. and Sunder, S. (1988). Rational Expectations and the Aggregation of Diverse Information in Laboratory Security Markets. \emph{Econometrica} 56(5), 1085--1118. \ddoi{10.2307/1911360}}
\dref{Rol1986a}{Roll, R. (1986). The Hubris Hypothesis of Corporate Takeovers. \emph{The Journal of Business} 59(2), 197--216. \ddoi{10.1086/296325}}
\dref{Smi1962a}{Smith, V.\,L. (1962). An Experimental Study of Competitive Market Behavior. \emph{Journal of Political Economy} 70(2), 111--137. \ddoi{10.1086/258609}}

\end{document}
